\documentclass[12pt,oneside]{article}

%%%%%%%%%%%%%%%%%%%%%%%%%%%%
%%   Zusaetzliche Pakete  %%
%%%%%%%%%%%%%%%%%%%%%%%%%%%%
\usepackage{enumerate}  
\usepackage{fancyhdr}
\usepackage{a4wide}
\usepackage{graphicx}
\usepackage{palatino}
\usepackage{multirow}
\usepackage{booktabs}
\usepackage{titlesec}
\usepackage{acronym}% http://ctan.org/pkg/acronym
\usepackage{enumitem}% http://ctan.org/pkg/enumitem

%folgende Zeile auskommentieren für englische Arbeiten
\usepackage[ngerman]{babel}
%folgende Zeile auskommentieren für deutsche Arbeiten
%\usepackage[ngerman, english]{babel}

\usepackage[T1]{fontenc}
\usepackage[utf8]{inputenc}
\usepackage[bookmarks]{hyperref}
\usepackage[justification=centering]{caption}
\usepackage[style=authoryear,natbib=true,backend=biber,maxbibnames=20]{biblatex}
\usepackage{csquotes}
\bibliography{literatur}

\setlength{\parindent}{0em} 
\setlist[itemize]{noitemsep, topsep=0pt}
\setlist[enumerate]{noitemsep, topsep=0pt}


%%%%%%%%%%%%%%%%%%%%%%%%%%%%%%
%% Definition der Kopfzeile %%
%%%%%%%%%%%%%%%%%%%%%%%%%%%%%%

\pagestyle{fancy}
\fancyhf{}
\cfoot{\thepage}
\setlength{\headheight}{16pt}

%%%%%%%%%%%%%%%%%%%%%%%%%%%%%%%%%%%%%%%%%%%%%%%%%%%%%
%%  Definition des Deckblattes und der Titelseite  %%
%%%%%%%%%%%%%%%%%%%%%%%%%%%%%%%%%%%%%%%%%%%%%%%%%%%%%

\newcommand{\JMUTitle}[9]{

  \thispagestyle{empty}
  \vspace*{\stretch{1}}
  {\parindent0cm
  \rule{\linewidth}{.7ex}}
  \begin{flushright}
    \vspace*{\stretch{1}}
    \sffamily\bfseries\Huge
    #1\\
    \vspace*{\stretch{1}}
    \sffamily\bfseries\large
    #2\\
    \vspace*{\stretch{1}}
    \sffamily\bfseries\small
    #3
  \end{flushright}
  \rule{\linewidth}{.7ex}

  \vspace*{\stretch{1}}
  \begin{center}
    \includegraphics[width=2in]{siegel} \\
    \vspace*{\stretch{1}}
    \Large Seminararbeit / Bachelorarbeit / Masterarbeit  \\

    \vspace*{\stretch{2}}
   \large Lehrstuhl f\"{u}r Wirtschaftsinformatik\\
    \large und Systementwicklung\\
    \large Universität Würzburg\\
    \vspace*{\stretch{1}}
    \large Betreuer:  #8 \\[1mm]
    
    \vspace*{\stretch{1}}
    \large W\"urzburg, den #7 \\
        \vspace*{\stretch{0.25}}

    Bearbeitungszeit: TT.MM.JJJJ - TT.MM.JJJJ % Die Bearbeitungszeit der Seminar-/ Abschlussarbeit ist hier einzutragen.

  \end{center}
}

\titlespacing*{\section}
{0pt}{3.5ex plus 1ex minus .2ex}{.2ex plus .2ex}
\titlespacing*{\subsection}
{0pt}{1.5ex plus 1ex minus .2ex}{.2ex plus .2ex}
\titlespacing*{\subsubsection}
{0pt}{1.5ex plus 1ex minus .2ex}{.2ex plus .2ex}




%%%%%%%%%%%%%%%%%%%%%%%%%%%%
%%  Beginn des Dokuments  %%
%%%%%%%%%%%%%%%%%%%%%%%%%%%%

\begin{document}

  \JMUTitle
      {TITEL}        % Titel der Arbeit
      {VORNAME NACHNAME}                        % Vor- und Nachname des Autors
      {MATRIKELNUMMER}
      
      {Wirtschaftswissenschaftlichen Fakultät}  % Name der Fakultaet
      {W"urzburg 2018}                          % Ort und Jahr der Erstellung
      {31.03.2018}                              % Tag der Abgabe
      {Prof. Dr. Frédéric Thiesse}               % Name des Erstgutachters
      {Zweitgutachter}                          % Name des Zweitgutachters

  \clearpage

\lhead{}
\pagenumbering{Roman} 
    \setcounter{page}{1}

\tableofcontents
\clearpage

%%%%%%%%%%%%%%%%%%%%%%%%%%%%
%%  Kurzzusammenfassung   %%
%%%%%%%%%%%%%%%%%%%%%%%%%%%%
\lhead{Zusammenfassung}
\section*{Zusammenfassung}
\addcontentsline{toc}{section}{Zusammenfassung}

Eine Kurzzusammenfassung der Vorgehensweise und der wesentlichen Ergebnisse.

Allgemeine Merkmale
\begin{itemize}
    \item Objektivität: Es soll sich jeder persönlichen Wertung enthalten.
    \item Kürze: Es soll so kurz wie möglich sein.
    \item Verständlichkeit: Es weist eine klare, nachvollziehbare Sprache und Struktur auf.
    \item Vollständigkeit: Alle wesentlichen Sachverhalte sollen enthalten sein.
    \item Genauigkeit: Es soll genau die Inhalte und die Meinung der Originalarbeit wiedergeben.
\end{itemize}{}

\newpage
\lhead{Abstract}
\section*{Abstract}
\addcontentsline{toc}{section}{Abstract}

Obige Zusammenfassung in englischer Sprache.

\newpage
\lhead{Abbildungsverzeichnis} % Bei englischsprachiger Arbeit anzupassen auf: List of Figures
\addcontentsline{toc}{section}{Abbildungsverzeichnis} % Bei englischsprachiger Arbeit anzupassen auf: List of Figures
\listoffigures

\newpage
\lhead{Tabellenverzeichnis} % Bei englischsprachiger Arbeit anzupassen auf: List of Tables
\addcontentsline{toc}{section}{Tabellenverzeichnis} % Bei englischsprachiger Arbeit anzupassen auf: List of Tables
\listoftables
\newpage

\setlength{\parskip}{0.5em} 


%%%%%%%%%%%%%%%%%%%%%%%%%%%%%%%%%%
%%  Definition der Abkürzungen  %%
%%%%%%%%%%%%%%%%%%%%%%%%%%%%%%%%%%
\lhead{Abkürzungsverzeichnis} % Bei englischsprachiger Arbeit anzupassen auf: List of Abbreviations
\section*{Abkürzungsverzeichnis} % Bei englischsprachiger Arbeit anzupassen auf: List of Abbreviations
\addcontentsline{toc}{section}{Abkürzungsverzeichnis} % Bei englischsprachiger Arbeit anzupassen auf: List of Abbreviations

\begin{acronym}
 \acro{ML}{Machine Learning}
\end{acronym}

%%%%%%%%%%%%%%%%%%%%%%%%%%%%
%%  Einstellungen  %%
%%%%%%%%%%%%%%%%%%%%%%%%%%%%
\clearpage
\pagenumbering{arabic}  
    \setcounter{page}{1}
\lhead{\nouppercase{\leftmark}}

%%%%%%%%%%%%%%%%%%%%%%%%%%%%
%%  Hauptteil  %%
%%%%%%%%%%%%%%%%%%%%%%%%%%%%

\section{Einleitung} \label{einleitung}
Dieser Teil der Arbeit sollte folgende Inhalte haben:

\begin{itemize}
\item Einführung in die Problemstellung
\item Motivation und Herleitung des Themas
\item Aufbau der Arbeit
\end{itemize}

Hinweis:
Es hat sich als hilfreich erwiesen, die Einleitung mit der Zusammenfassung bzw. dem  Abstract und der Schlussfolgerung zu vergleichen. Damit stellt man sicher, dass diese inhaltlich im Bezug auf Zielsetzung und Motivation übereinstimmen. Der Umfang sollte ca. 5 \% der gesamten Arbeit betragen.


\section{Abschnitt} 
Hier soll eine kurze Einführung erfolgen, die den Zusammenhang des Kapitels zur Arbeit herstellt.

Generelle Hinweise:

\begin{itemize}
\item eindeutige Begrifflichkeit verwenden
\item auf logische Herleitung der Argumentation achten
\end{itemize}

Ein Hauptabschnitt - idealerweise sollten keine Abschnitte leer sein.

\subsection{Unterabschnitt}
Ein Unterabschnitt - idealerweise sollten keine Abschnitte leer sein.

\subsubsection{Unterunterabschnitt}
Ein Unterunterabschnitt - idealerweise sollten keine Abschnitte leer sein.

\section{Einfache Formatvorlagen}

\textbf{Das ist fett gedruckter Text}.

\textit{Das ist kursiver Text}.


Auflistungen sind oft hilfreich für die Strukturierung:
\begin{itemize}
    \item Erster Eintrag
    \item Zweiter Eintrag
\end{itemize}

Nummerierte Aufzählungen sind oft hilfreich für Reihenfolgen:
\begin{enumerate}
    \item Erster Eintrag
    \item Zweiter Eintrag
\end{enumerate}


\section{Zitieren und Referenzieren}

Beiträge in Fachzeitschriften wie \citet[12]{clemen1989combining} oder Konferenzartikel wie \citet[6]{he2017mask} werden auf diese Weise im Text zitiert. In anderen Fällen möchte man aber in Klammern zitieren \citep[10]{clemen1989combining}, auch mit mehreren Autoren (\citealp[3]{baumol1958warehouse}; \citealp[15]{clemen1989combining}; \citealp[12]{he2017mask}).

Die Seitenzahlen müssen angegeben werden \citep[28]{chollet2018deep}. Bezieht sich das Zitat auf eine Textstelle, die sich über mehrere Seiten streckt, so sind diese entsprechend anzugeben: \citet[28-29]{chollet2018deep} bzw. \citet[28-35]{chollet2018deep}

So wird eine Webquelle zitiert: \citet{shiny1}. Es kann bei kurzen Informationen im Internet aber auch reichen die Adresse\footnote{\url{https://shiny.rstudio.com/tutorial/written-tutorial/lesson1/}} als Fußnote einzubetten.

Bei einem direkten Zitat muss der zitierte Text originalgetreu wiedergegeben werden. Rechtschreibfehler oder eine veraltete Orthographie werden unverändert wiedergegeben. „Der zitierte Text steht immer in Anführungszeichen"\ \citep[28]{chollet2018deep}. 

So werden andere Teile der Arbeit referenziert: Kapitel \ref{einleitung}, Gleichung \ref{eq:1} zeigt...

So verweisen wir auf eine Fußnote \footnote{dies ist eine Fußnote}.

\section{Abbildungen}

Abbildungen erfordern das package \textit{graphicx}. 
Idealerweise verwendet man Vektorgrafiken oder hochaufgelöste Bitmaps. 
Eine gute Variante ist das Verwenden von PDFs.

\begin{figure}[h]
    \centering
    \includegraphics[width=0.3\textwidth]{siegel.pdf}
    \caption{Siegel der Universität}
    \label{fig:my_label}
\end{figure}


\section{Tabellen}

Die Tabular-Umgebung gibt die Anzahl Spalten an, deren Orientierung, Breite und evtl. Zwischenlinien. 


\begin{table}[ht]
    \centering
    \caption{Meine Tabelle}
        \begin{tabular}{ cccc } 
        \toprule
        col1 & col2 & col3 \\
        \midrule
        \multirow{3}{4em}{Multiple row} & cell2 & cell3 \\ 
        & cell5 & cell6 \\ 
        & cell8 & cell9 \\ 
        \bottomrule
    \end{tabular}
    \label{tab:countries}
\end{table}

\section{Formeln}

\begin{equation}
    \sum_{i=1}^N x_i
    \label{eq:1}
\end{equation}

\section{Abkürzungen}
Bei der ersten Verwendung wird die Abkürzung eines Fachbegriffs wie zum Beispiel \ac{ML} eingeführt und daher ausgeschrieben. Bei der zweiten Verwendung der Abkürzung \ac{ML} ist dies nicht mehr nötig. Die Abkürzungen sind in dem Abschnitt \textit{Definition der Abkürzungen} einzupflegen. Das Abkürzungsverzeichnis ist alphabetisch anzuordnen. Bleibt das Abkürzungsverzeichnis leer, so kann dieser Abschnitt (Zeilen 164-173 im main.tex) gelöscht werden.


\section{Fazit}
In der Schlussfolgerung sollen

\begin{itemize}
\item die Themenstellung
\item der gewählte Ansatz
\item die Ergebnisse der Arbeit
\item eine kritische Stellungnahme/Einschätzung
\item nächste Schritte
\end{itemize}
deutlich werden.

Hinweis:
Die Schlussfolgerung sollte mit der Zusammenfassung bzw. dem Abstract und der Einleitung abgeglichen werden. Es sollte immer eine Zusammenfassung der wesentlichen Er-kenntnisse der eigenen Arbeit sein, die den Forschungsbeitrag darstellt. Der Umfang der Schlussfolgerung sollte ähnlich wie die Einleitung ca. 5 \% der gesamten Arbeit betragen.


%%%%%%%%%%%%%%%%%%%%%%%%%%%%
%% Literaturverzeichnis wird 
%% automatisch eingefügt
%%%%%%%%%%%%%%%%%%%%%%%%%%%%
\clearpage
\lhead{}
\printbibliography
\addcontentsline{toc}{section}{\bibname}


%%%%%%%%%%%%%%%%%%%%%%%%%%%%
%% Anhang (optional) 
%%%%%%%%%%%%%%%%%%%%%%%%%%%%
\clearpage
\appendix
\section{Anhang A} % Bei englischsprachiger Arbeit anzupassen auf: Appendix A

%%%%%%%%%%%%%%%%%%%%%%%%%%%%
%% Versicherung zur Leistungserbringung 
%% ist auf der Webseite des 
%% Lehrstuhls zu finden.
%%%%%%%%%%%%%%%%%%%%%%%%%%%%

\end{document}
